% This is part of Mes notes de mathématique
% Copyright (c) 2011-2018
%   Laurent Claessens, Carlotta Donadello
% See the file fdl-1.3.txt for copying conditions.

%+++++++++++++++++++++++++++++++++++++++++++++++++++++++++++++++++++++++++++++++++++++++++++++++++++++++++++++++++++++++++++
\section{Espaces métriques}
%+++++++++++++++++++++++++++++++++++++++++++++++++++++++++++++++++++++++++++++++++++++++++++++++++++++++++++++++++++++++++++

%---------------------------------------------------------------------------------------------------------------------------
\subsection{Espaces métrisables}
%---------------------------------------------------------------------------------------------------------------------------

\begin{definition}
    Un espace topologique est \defe{métrisable}{espace!topologique!métrisable} s'il est homéomorphe à un espace métrique.
\end{definition}


\begin{definition}		\label{DefCvSuiteEGVN}
	Soit une suite $(x_n)$ dans un espace vectoriel normé $V$. Nous disons qu'elle est
    \defe{convergente}{convergence!dans un espace vectoriel normé} s'il existe un élément $\ell\in V$ tel que
	\begin{equation}
		\forall \varepsilon>0,\,\exists N\in\eN\tq n\geq N\Rightarrow \| x_n-l \|<\varepsilon.
	\end{equation}
	Dans ce cas, $\ell$ est appelé la \defe{limite}{limite!suite} de la suite $(x_n)$.
\end{definition}


\begin{proposition} \label{PROPooKNVUooMbLZoy}
    Une fonction séquentiellement continue sur un espace métrisable et à valeurs dans un espace métrique est continue.
\end{proposition}

\begin{proof}
    Soient \( E\) un espace métrique et \( \phi\colon X\to (E,d)\) un homéomorphisme. Nous supposons que \( f\colon X\to Y\) est séquentiellement continue. Nous considérons l'application \( \tilde f=f\circ\phi^{-1}\), c'est à dire
    \begin{equation}
        \begin{aligned}
            \tilde f\colon E&\to Y \\
            a&\mapsto f\big( \phi^{-1}(a) \big).
        \end{aligned}
    \end{equation}
    L'application \( \phi^{-1}\) est continue et donc séquentiellement continue. De plus \( \tilde f\) est séquentiellement continue. En effet si \( a_k\stackrel{E}{\longrightarrow}a\), alors
    \begin{equation}
        \tilde f(a_k)=f\big( \phi^{-1}(a_k) \big),
    \end{equation}
    mais \( \phi^{-1}\) est séquentiellement continue, donc \( \phi^{-1}(a_k)\stackrel{X}{\longrightarrow}\phi^{-1}(a)\), ce qui signifie que \( \phi^{-1}(a_k)\) est une suite convergente dans \( X\) et donc
    \begin{equation}
        \lim_{k\to \infty} \tilde f(a_k)=\lim_{k\to \infty} f\big( \phi^{-1}(a_k) \big)=f\big( \phi^{-1}(a) \big)=\tilde f(a).
    \end{equation}
    L'application \( \tilde f\) est donc séquentiellement continue. Mais étant donné que \( \tilde f\) est définie sur un espace métrique (\( E\)) et à valeurs dans un métrique, elle est continue par la proposition~\ref{PropXIAQSXr}. L'application \( f=\tilde f\circ\phi\) est donc continue en tant que composée d'applications continues.
\end{proof}

%---------------------------------------------------------------------------------------------------------------------------
\subsection{Fonctions continues}
%---------------------------------------------------------------------------------------------------------------------------

\begin{theorem}[Théorème \wikipedia{fr}{Théorème_des_fermés_emboités}{des fermés emboîtés}\cite{OIywOjl}]   \label{ThoCQAcZxX}
    Soit \( (E,d)\) un espace métrique. Il est complet si et seulement si toute suite décroissante de fermés non vides dont le diamètre tend vers zéro a une intersection qui se réduit à un seul point.
\end{theorem}

\begin{proof}
    \begin{subproof}
    \item[Condition suffisante]

        Soit \( \{ F_n \}_{n\in \eN}\) une telle suite de fermés emboités. Si nous choisissons des points \( x_n\in F_n\), nous obtenons une suite \( (x_n)\) de Cauchy et qui est par conséquent convergente vu que l'espace est par hypothèse complet. De plus, pour chaque \( N\geq n\), la queue de suite \( (x_n)_{n\geq N}\) est contenue dans \( F_N\) et donc converge vers un élément de \( F_N\) (parce que ce dernier est fermé). Donc la limite de \( (x_n)\) est dans \( \bigcap_{n\in \eN}F_n\).

        De plus cette intersection a diamètre nul parce que le diamètre de \( \bigcap_{n\in \eN}F_n\) est majoré par tous les diamètres des \( F_n\), lesquels sont arbitrairement petits par hypothèse. Donc l'intersection est réduite a un point.

    \item[Condition nécessaire]

        Soit \( (x_n)\) un suite de Cauchy. Nous considérons les ensembles
        \begin{equation}
            F_n=\overline{ \{ x_i\tq i\geq n \} }.
        \end{equation}
        Le fait que la suite soit de Cauchy implique que \( \diam(F_n)\to 0\). Par hypothèse, nous avons alors
        \begin{equation}
            \bigcap_{n\in \eN}F_n=\{ a \}.
        \end{equation}
        Pour s'assurer que \( a\) est bien la limite de \( (x_n)\), il suffit de remarquer que
        \begin{equation}
            d(x_n,a)\leq \diam F_n\to 0.
        \end{equation}
    \end{subproof}
\end{proof}


\begin{lemma}       \label{LemooynkH}
    Soit \( A_n\) une suite décroissante de fermés dans un espace métrique\footnote{L'hypothèse métrique provient de l'utilisation de Bolzano-Weierstrass, lequel est vrai pour les espaces séquentiellement compacts, dont les espaces métriques.} compact \( K\). Alors
    \begin{equation}
        C=\bigcap_{n\in \eN}A_n
    \end{equation}
    est non vide.
\end{lemma}

\begin{proof}
    Soit \( (x_n)\) une suite dans \( K\) telle que \( x_n\in A_n\). La suite étant contenue dans \( A_1\), et \( A_1\) étant compact (lemme~\ref{LemnAeACf}), elle possède une sous-suite \( (y_n=x_{\sigma_1(n)})\) convergente dont la limite est dans \( A_1\) par le théorème de Bolzano-Weierstrass~\ref{ThoBWFTXAZNH}. Une queue de la suite \( y_n\) est dans \( A_2\) et nous considérons donc une sous-suite convergente dans \( A_2\) donnée par
    \begin{equation}
        z_n=y_{\sigma_2(n)}=x_{\sigma_1\sigma_2(n)}.
    \end{equation}
    En continuant ainsi nous construisons une suite convergente dans \( A_k\). Nous considérons enfin la suite
    \begin{equation}
        y_n=x_{\sigma_1\ldots \sigma_n(n)}.
    \end{equation}
    Pour tout \( k\), une queue de cette suite est une sous-suite de \( x_{\sigma_1\ldots \sigma_k(n)}\) et par conséquent cette suite converge dans \( A_k\). La limite de cette suite est donc dans l'intersection demandée.
\end{proof}

\begin{remark}
    Cette propriété est fausse pour les ouverts. Par exemple
    \begin{equation}
        \bigcap_{n>1}\mathopen] 0 , \frac{1}{ n } \mathclose[=\emptyset.
    \end{equation}
\end{remark}

\begin{lemma}   \label{LemKIcAbic}
    Si \( K\) est un compact dans un espace métrique et \( F\) un fermé disjoint de \( K\), alors \( d(K,F)>0\).
\end{lemma}

\begin{proof}
    La fonction
    \begin{equation}
        \begin{aligned}
             K&\to \eR \\
            x&\mapsto d(x,F)
        \end{aligned}
    \end{equation}
    est une fonction continue sur \( K\), et donc atteint son minimum par le théorème de Weierstrass~\ref{ThoWeirstrassRn}. Soit \( x_0\in K\) un point de \( K\) qui réalise ce minimum. Si \( d(x_0,F)=0\), alors on aurait une suite \( (x_n)\) dans \( F\) qui convergerait vers \( x_0\), mais \( F\) étant fermé cela signifierait que \( x_0\) serait dans \( F\), ce qui contredirait l'hypothèse que \( F\) et \( K\) sont disjoints.
\end{proof}

\begin{proposition}[\cite{AntoniniAndAl-EspacesMetriquesCompacts}]
    Une isométrie d'un espace métrique compact sur lui-même est une bijection.
\end{proposition}

\begin{proof}
    Soient \( X\) un espace métrique compact et \( f\colon X\to X\) une isométrie. Le fait que \( f\) soit injective est obligatoire (sinon il y a des images dont la distance est nulle). Il faut montrer que \( f\) est surjective.

    Soit \( x\in X\) hors de \( f(X)\). Le lemme~\ref{LemKIcAbic} appliqué au fermé \( \{ x \}\) et au compact \( f(K)\) donne un \( r>0\) tel que
    \begin{equation}
        d\big( x,f(K)\big)>r.
    \end{equation}
    Soit la suite \( u_n=f^n(x)\); c'est une suite dans \( K\) et possède donc une sous-suite convergente (Bolzano-Weierstrass\ref{ThoBWFTXAZNH}) que l'on nomme \( (y_n)\). Vu que \( f\) est une isométrie,
    \begin{equation}
        d(y_{n},y_{n+1})=d(x,y_m)>r
    \end{equation}
    pour un certain \( m\leq n+1\). Cela signifie que pour tout \( n\), nous avons \( d(y_n,y_{n+1})>r\), ce qui contredit le fait que la suite \( (y_n)\) converge.
\end{proof}

\begin{proposition} \label{PropLHWACDU}
    Soient \( (X,d)\) un espace métrique compact et \( (u_n)\) une suite de \( X\) telle que
    \begin{equation}
        \lim_{n\to \infty} d(u_n,u_{n+1})=0.
    \end{equation}
    Alors l'ensemble des points d'accumulation\footnote{Définition \ref{DEFooGHUUooZKTJRi}.} de \( (u_n)\) est connexe.
\end{proposition}
\index{connexité!points d'accumulation}
\index{compacité}

\begin{proof}
    Nous notons \( \Gamma\) l'ensemble des points d'accumulation de la suite.
    \begin{subproof}
    \item[\( \Gamma\) est compact]
        Nous notons \( A_p=\{ u_n\tq n\geq p \}\) et nous avons
        \begin{equation}
            \Gamma=\bigcap_{p\in \eN}\overline{ A_p }
        \end{equation}
        parce que si \( x\in\Gamma\), alors pour tout \( n\), il existe \( m>n\) tel que \( x_m\in B(x,\epsilon)\), et donc tel que \( x\in B(x_m,\epsilon)\). Donc pour tout \( \epsilon\) et pour tout \( p\), l'intersection \( B(x,\epsilon)\cap A_p\) est non vide.

        En tant qu'intersection de fermés, \( \Gamma\) est fermé (lemme~\ref{LemQYUJwPC}). En tant que fermé dans un compact, \( \Gamma\) est compact (lemme~\ref{LemnAeACf}).

    \item[Recouvrement par deux compacts]

        Supposons que \( \Gamma\) ne soit\quext{est-ce qu'il faut vraiment un subjonctif ici ?} pas connexe. Nous pouvons alors considérer \( S\) et \( O\), deux ouverts disjoints recouvrant \( \Gamma\) et intersectant tout deux \( \Gamma\). Nous posons alors
        \begin{subequations}
            \begin{align}
                A&=S\cap\Gamma\\
                B&=O\cap\Gamma,
            \end{align}
        \end{subequations}
        et nous avons évidemment \( \Gamma=A\cup B\). Montrons que \( A\) est fermé (\( B\) le sera aussi par le même raisonnement). Soit une suite d'éléments de \( S\cap \Gamma\) convergent dans \( X\). Alors la limite est dans \( \bar\Gamma=\Gamma\) et donc elle est donc \( O\) ou \( S\), mais elle est certainement dans \( \bar S\). Cependant \( \bar S\) n'intersecte pas \( O\). En effet si \( x\in \bar S\cap O\), alors tout voisinage de \( x\) intersecterait \( S\), mais il y a des voisinages de \( x\) étant inclus dans \( O\) parce que \( O\) est ouvert; cela donnerait une intersection entre \( O\) et \( S\), ce qui est impossible. Donc la limite n'est pas dans \( O\) et donc elle est dans \( S\). Au final la limite est dans \( S\cap \Gamma\), ce qui prouve son caractère fermé.

        Comme d'habitude, \( \Gamma\cap S\) est compact parce que fermé dans un compact.

    \item[Décomposition en trois morceaux]

        Vu que \( A\) et \( B\) sont des compacts disjoints, nous avons \( d(A,B)=\alpha>0\) pour un certain \( \alpha\) par le lemme~\ref{LemKIcAbic}. Nous notons
        \begin{subequations}
            \begin{align}
                A'&=\{x\in X\tq d(x,A)<\frac{ \alpha }{ 3 }\}\\
                B'&=\{x\in X\tq d(x,B)<\frac{ \alpha }{ 3 }\}
            \end{align}
        \end{subequations}
        Nous avons \( A'=\bigcup_{x\in A}B(x,\frac{ \alpha }{ 3 })\) et donc en tant qu'union d'ouverts, \( A'\) est ouvert (définition de la topologie). Même chose pour \( B'\).

        Enfin nous notons
        \begin{equation}
            K=X\setminus(A'\cup B')
        \end{equation}
        qui est fermé en tant que complémentaire d'ouvert, et donc compact. Étant donné que \( A\subset A'\) et \( B\subset B' \), nous avons \( K\cap \Gamma=\emptyset\).

        L'idée est maintenant de montrer que \( K\) contient un point d'accumulation de \( (u_n)\).

    \item[Sous-suites de \( (u_n)\)]

        L'hypothèse sur la suite \( (u_n)\) nous indique qu'il existe un \( N_0\) tel que \( \forall n\geq N_0\),
        \begin{equation}    \label{EqIHioHjW}
            d(u_{n},u_{n+1})<\frac{ \alpha }{ 3 }.
        \end{equation}
        Soient \( N>N_0 \) et \( x_0\in A\). Étant donné que \( x_0\) est point d'accumulation de la suite, il existe \( n_1>N\) tel que \( d(x_0,u_{n_1})<\frac{ \alpha }{ 3 }\). Même chose dans \( B\) : nous prenons \( y_0\in B\) et un naturel \( n_2>n_1\) tel que \( d(y_0,u_{n_2})<\frac{ \alpha }{ 3 }\). Nous avons \( u_{n_1}\in A'\) et \( u_{n_2}\in B'\).

        Soit \( n_0\) le plus petit naturel supérieur à \( n_1\) tel que \( u_{n_0}\notin A'\). Cela existe parce que \( u_{n_2}\in B'\) et \( B'\cap A'=\emptyset\), mais \( n_0\) n'est pas \( n_2\) lui-même parce que \( d(A',B')\geq \frac{ \alpha }{ 3 }\) alors que nous considérons \( n_0,n_1,n_2>N_0\) et donc pour tous les \( i\) entre \( n_1\) et \( n_2\) (compris), \( d(u_i,u_{i+1})<\frac{ \alpha }{ 3 }\). Notons qu'ici le strict dans la condition \eqref{EqIHioHjW} est important. Nous avons donc \(N_0<n_1<n_0<n_2\).

        Nous allons maintenant montrer que \( u_{n_0}\) est dans \( K\). C'est fait pour : il est loin en même temps de \( A'\) et de \( B'\). En utilisant l'inégalité triangulaire à l'envers, nous avons
        \begin{equation}
            \begin{aligned}[]
            d(u_{n_0},B)&\geq d(u_{n_0-1},B)-d(u_{n_0-1},u_{n0})\\
            &\geq d(A,B)-d(u_{n_0-1},A)-d(u_{n_0-1},u_{n_0})\\
            &\geq \alpha-\frac{ \alpha }{ 3 }-\frac{ \alpha }{ 3 }\\
            &=\frac{ \alpha }{ 3 }.
            \end{aligned}
        \end{equation}
        Pour la dernière inégalité nous avons utilisé le fait que \( u_{n_0-1}\) n'est pas dans \( A'\). Bref, nous avons montré que \( u_{n_0}\) n'est pas dans \( B'\) (dans la définition de ce dernier nous avons bien une inégalité stricte). Vu que par définition \( u_{n_0}\) n'est pas non plus dans \( A'\), nous avons \( u_{n_0}\in K\).

        Nous avons montré jusqu'à présent que pour tout \( N\geq N_0\), il existe un \( n_0\geq N\) tel que \( u_{n_0}\in K\). Cela nous construit donc une sous-suite \( (v_n)\) de \( (u_n)\) contenue dans \( K\). En tant que suite dans le compact \( K\), la suite \( (v_n)\) admet un point d'accumulation dans \( K\). Ce point est également point d'accumulation de la suite \( (u_n)\) complète, ce qui donne un point d'accumulation de \( (u_n)\) dans \( K\) et donc une contradiction.

    \end{subproof}
    Nous concluons que \( \Gamma\) est connexe.
\end{proof}

Encore une petite conséquence sans ambition du théorème de Bolzano-Weierstrass.
\begin{proposition}\label{PropHNylIAW}
    Si \( (x_n)\) est une suite dans un compact telle que toute sous-suite convergente ait le même point \( x\) comme limite. Alors la suite entière converge vers \( x\).
\end{proposition}

\begin{proof}
    Supposons que ce ne soit pas le cas. Alors il existe un \( \epsilon\) tel que pour tout \( N>0\), il existe \( n>N\) avec \( d(x_n,x)>\epsilon\). Cela nous donne une sous-suite de \( (x_n)\) composée d'éléments tous à une distance de \( x\) supérieure à \( \epsilon\). Nous la nommons \( (y_n)\); c'est une suite dans un compact qui admet donc une sous-suite convergente (et une telle sous-suite est une sous-suite de \( (x_n)\)) dont la limite devrait être \( x\), mais c'est impossible par construction.
\end{proof}

\begin{lemmaDef}       \label{LemGDeZlOo}
    Soi \( \Omega\) un ouvert dans un espace métrique \( E\). Il existe une suite \( (K_n)\) de compacts tels que
    \begin{enumerate}
        \item
            \( K_n\subset \Omega\)
        \item
            \( \bigcup_{n=0}^{\infty}K_n=\Omega\)
        \item
            \( K_n\subset\Int(K_{n+1})\).
    \end{enumerate}
    Une telle suite de compacts vérifie alors
    \begin{enumerate}
        \item
            Il existe \( \delta_n\) tel que pour tout \( z\in K_n\), \( B(z,\delta_n)\subset K_{n+1}\).
        \item
            Tout compact de \( \Omega\) est inclus dans \( \Int(K_n)\) pour un certain \( n\).
    \end{enumerate}
    Une telle suite de compacts est une \defe{suite exhaustive}{compact!suite exhaustive}\index{exhaustive (suite de compacts)!} de compacts pour \( \Omega\).
\end{lemmaDef}

\begin{proof}
    Nous considérons les ensembles
    \begin{equation}
        V_n=\{ z\in E\tq | z | \}\cup\bigcup_{a\notin\Omega}B(a,\frac{1}{ n }),
    \end{equation}
    et nous définissons \( K_n=\complement V_n\). Vérifions que ces ensembles vérifient tout ce qu'il faut.
    \begin{enumerate}
        \item
            Si \( a\notin\Omega\) alors \( a\) est dans tous les \( V_n\) et donc dans aucun des \( K_n\); nous avons donc bien \( K_n\subset\Omega\).
        \item
            Si \( z\in \Omega\) alors nous prenons \( n_1>| z |\) puis \( n_2\) tel que \( B(z,\frac{1}{ n_2 })\subset \Omega\). Alors \( z\in K_n\) avec \( n>\max(n_1,n_2)\).
        \item
            Une chose à comprendre est que si \( z\in K_n\), alors \( d(z,\complement \Omega)\geq \frac{1}{ n }\). Du coup si nous prenons \( \delta\) tel que
            \begin{equation}
                \frac{1}{ n+1 }<\delta<\frac{1}{ n }
            \end{equation}
            alors \( B(z,\delta)\subset K_{n+1}\).
        \item
            Enfin, les \( K_n\) sont tous compacts. En effet ils sont bornés parce que \( K_n\subset B(0,n)\) et ensuite \( K_n\) est fermé en tant que complémentaire d'un ouvert (\( V_n\) est ouvert en tant qu'union d'ouverts).
    \end{enumerate}

    Nous passons maintenant aux propriétés, qui sont indépendantes de la façon dont nous avons construit les \( K_n\) vérifiant les conditions.
    \begin{enumerate}
        \item

            Nous pouvons considérer la fonction \( K_n\to \eR\) donnée par \( z\mapsto d(z,\complement K_{n+1})\). Vu que \( K_n\subset\Int(K_{n+1})\), c'est une fonction (continue sur le compact \( K_n\)) prenant des valeurs strictement positives. Elle a donc un minimum strictement positif. Si \( \delta_n\) est plus petit que ce minimum nous avons \( B(z,\delta_n)\subset K_{n+1}\) pour tout \( z\in K_n\).

        \item

            D'abord nous avons \( \Omega=\bigcup_{n=0}^{\infty}\Int(K_n)\). En effet nous avons
            \begin{equation}
                \Omega=\bigcup_{n=0}^{\infty}K_n\subset\bigcup_{n=0}^{\infty}\Int(K_{n+1})\subset\bigcup_{n=0}^{\infty}\Int(K_n).
            \end{equation}
            L'inclusion dans l'autre sens est facile.

            Soit \( K\) compact dans \( \Omega\). Vu que \( \Omega\) est l'union des \( \Int(K_n)\), nous avons
            \begin{equation}
                K\subset\bigcup_{n=0}^{\infty}\Int(K_n).
            \end{equation}
            Cela donne à \( K\) un recouvrement par des ouverts dont nous pouvons extraire un sous-recouvrement fini par compacité. Les \( K_n\) étant croissants, du recouvrement fini, il suffit de prendre le plus grand (disons \( K_m\)) et nous avons \( K\subset\Int(K_m)\).

    \end{enumerate}
\end{proof}
Notons qu'avec la suite de \( K_n\) telle que construite, le dernier point est réglé en prenant
\begin{equation}
    \frac{1}{ n+1 }<\delta_n<\frac{1}{ n }.
\end{equation}


\begin{theorem}[Tykhonov]\index{théorème!Tykhonov}\label{ThoFWXsQOZ}
    Un produit quelconque d'espaces métriques non vides est compact si et seulement si chacun de ses facteurs est compact.
\end{theorem}
Nous n'allons donner la preuve que dans le cas d'un produit fini dans le théorème~\ref{THOIYmxXuu}.

%---------------------------------------------------------------------------------------------------------------------------
\subsection{Ensembles enchaînés}
%---------------------------------------------------------------------------------------------------------------------------

Soit \( (x,d)\) un espace métrique.
\begin{definition}
    Une \defe{\( \epsilon\)-chaîne}{chaîne} joignant les points \( a\) et \( b\) de \( X\) est une suite finie \( (u_0,\ldots, u_n)\) dans \( X\) telle que \( u_0=a\), \( u_n=b\) et pour tout \( 0\leq i\leq n-1\) nous avons \( d(u_n,u_{n+1})\leq \epsilon\).

    Une partie \( A\) de \( X\) est \defe{bien enchaînée}{bien!enchaîné} si pour tout \( \epsilon>0\) et pour tout \( a,b\in A\), il existe une \( \epsilon\)-chaîne joignant \( a\) et \( b\) dans $A$.
\end{definition}
Les rationnels dans \( \eR\) sont bien enchaînés.

\begin{proposition}
    Un espace connexe est bien enchaîné.
\end{proposition}
%TODO: une preuve.

\begin{proposition}
    La fermeture d'un ensemble bien enchaîné dans un espace métrique compact \( (X,d)\) est connexe.
\end{proposition}
\index{connexité}
\index{compacité}

\begin{proof}
    Soit \( A\subset X\) un ensemble bien enchaîné, et soient \( a,b\in \bar A\). Nous construisons une suite \( (u_k)\) dans \( A\) de la façon suivante. Pour chaque \( n>0\) nous prenons \( a'\in B(a,\frac{1}{ n })\cap A\) et \( b'\in B(b,\frac{1}{ n })\cap A\). Ensuite nous considérons une \( \frac{1}{ n }\)-chaîne \( \{ v_i^{(n)} \}_{i\in I_n}\) dans \( A\) entre \( a'\) et \( b'\). Ici l'ensemble \( I_n\) est fini. La suite \( (u_k)\) est simplement construite en mettant bout à bout les éléments \( v_i^{(n)}\).

    La suite ainsi construite est une suite dans \( A\) admettant \( a\) et \( b\) comme points d'accumulation (les autres points d'accumulation sont également dans \( \bar A\)) et telle que \( \lim_{k\to \infty} d(u_k,u_{k+1})=0\). Par conséquent la proposition~\ref{PropLHWACDU} nous dit que l'ensemble des points d'accumulation de \( (u_k)\) est connexe dans \( X\). Nous le notons \( C_{a,b}\).

    Si nous fixons \( a\in \bar A\), alors nous avons
    \begin{equation}
        \bigcup_{x\in \bar A}C_{a,x}=\bar A.
    \end{equation}
    Vu que le membre de gauche est une union de connexes, c'est un connexe par la proposition~\ref{PropIWIDzzH}.
\end{proof}

\begin{corollary}       \label{CORooSIKCooTncoQm}
    Un espace métrique compact est connexe si et seulement s'il est bien enchaîné.
\end{corollary}

%---------------------------------------------------------------------------------------------------------------------------
\subsection{Produit fini d'espaces métriques}
%---------------------------------------------------------------------------------------------------------------------------

\begin{definition}\label{DefZTHxrHA}
    Si \( (E_1,d_1)\),\ldots, \( (E_n,d_n)\) sont des espaces métriques nous mettons la distance suivante sur le produit cartésien \( E=E_1\times\ldots\times E_n\) :
    \begin{equation}
        d(x,y)=\max_{i=1,\ldots, n}d_i(x_i,y_i).
    \end{equation}
\end{definition}

\begin{theorem}[\cite{MonCerveau}]\label{THOIYmxXuu}
    Un produit fini d'espaces métriques non vides est compact si et seulement si chacun de ses facteurs est compact.
\end{theorem}
\index{compact!produit fini}
\index{théorème!Tykhonov!fini}

\begin{proof}
    Soient \( K_1\),\ldots, \( K_n\) des compacts et \( K=K_1\times \ldots\times K_n\) le produit muni de sa métrique usuelle de la définition \eqref{DefZTHxrHA} (attention : chacun des \( K_i\) peut être de dimension infinie) :
    \begin{equation}
        d(\alpha,\beta)=\max\{ d_i(\alpha_i,\beta_i) \}
    \end{equation}
    où \( d_i\) est la distance sur \( K_i\). Si \( (\alpha_n)\) est une suite dans \( K\) alors la suite \( (\alpha_n)_1\) est une suite dans le compact \( K_1\) dont nous pouvons extraire une sous-suite convergente (Bolzano-Weierstrass~\ref{ThoBWFTXAZNH}). De la sous-suite de \( \alpha\) correspondante nous extrayons la sous-suite pour la seconde composante, etc.

    En fin de compte nous avons une sous-suite (que nous nommons \( \alpha\) également) donc chacune des composantes est convergente. Nous nommons \( \ell_k\) les limites correspondantes. Soit \( \epsilon>0\) pour chaque \( k=1,\ldots, n\), il existe \( N_k>0\) tel que si \( p>N_k\) alors
    \begin{equation}
        d\big( (\alpha_p)_k-\ell_p \big)\leq \epsilon.
    \end{equation}
    Ici \( \alpha_p\in K\) est le \( p\)\ieme élément de la suite \( \alpha\) et \( (\alpha_p)_i\in K_i\) est la \( i\)\ieme composante de \( \alpha_p\). En prenant \( N=\max_kN_k\) et \( n>N\) nous avons
    \begin{equation}
        d\big( \alpha_n,(\ell_1,\ldots, \ell_n) \big)\leq\epsilon.
    \end{equation}
    Par conséquent de la suite \( (\alpha)\) nous avons extrait une sous-suite convergente et la partie «réciproque» de Bolzano-Weierstrass nous assure alors que \( K\) est compact.

    À l'inverse si un des facteurs n'est pas compact (mettons \( K_1\)) alors nous prenons un recouvrement \( \{ \mO_i \}_{i\in I}\) de \( K_1\) par des ouverts duquel il est impossible d'extraire un sous-recouvrement fini. Ensuite nous posons
    \begin{equation}
        \mP_i=\mO_i\times K_2\times\ldots\times K_n,
    \end{equation}
    qui est un recouvrement de \( K\) par des ouverts (de \( K\)) d'où aucun sous-recouvrement fini ne peut être extrait.
\end{proof}

Pour la culture générale, il y a bien entendu moyen de faire des produits dénombrables et pire d'espaces métriques.
\begin{definition}[\cite{AntoniniAndAl-TheoremeTykhonov}]
    Soient \( (E_n,d_n)\) des espaces métriques. Sur l'ensemble produit \( E=\prod_{i=1}^{\infty}E_i\) nous définissons la métrique
    \begin{equation}
        d(x,y)=\sum_{k=1}^{\infty}\frac{1}{ 2^k }d'_k(x_i,y_i)
    \end{equation}
    où \( d'_i=\min(d_i,1)\).
\end{definition}
On peut montrer que ce \( d\) est bien une distance et que \( (E,d)\) devient un espace métrique.

\begin{theorem}[Tykhonov dénombrable\cite{AntoniniAndAl-TheoremeTykhonov}] \label{ThoKKBooNaZgoO}  % Ce résultat n'est pas censé être utilisé dans l'agrégation.
    Un produit dénombrable d'espaces métriques non vides est compact si et seulement si chacun de ses facteurs est compact.
\end{theorem}
\index{compact!produit dénombrable}
\index{théorème!Tykhonov!dénombrable}
Note : ce résultat est encore valable pour un produit quelconque, c'est le théorème de Tykhonov~\ref{ThoFWXsQOZ}.

%--------------------------------------------------------------------------------------------------------------------------- 
\subsection{Équicontinuité}
%---------------------------------------------------------------------------------------------------------------------------

\begin{definition}[\cite{ooYDVWooPWLUGW}]       \label{DEFooSGMVooASNbxo}
    Soit une famille de fonctions \( f_i\colon X\to E\) indexée par un ensemble \( I\) où \( X\) est un espace topologique et \( E\) un espace métrique. Cette famille est \defe{équicontinue}{famille équicontinue} en \( x\in X\) si pour tout \( \epsilon>0\), il existe un voisinage \( V\) de \( x\) tel que
    \begin{equation}
        \| f_i(x)-f_i(y) \|<\epsilon
    \end{equation}
    pour tout \( i\) dès que \( x,y\in V\).
\end{definition}

La proposition suivante permet de montrer que certaines fonctions définies par une limite sont continues. Ce sera par exemple le cas de la fonction puissance, proposition \ref{PROPooUQNZooSSHLqr}.
\begin{proposition}[\cite{ooYDVWooPWLUGW}]     \label{PROPooICNNooAMjcut}
    Si \( \{ f_i \}\) est une suite équicontinue de fonctions, et si il y a convergence simple \( f_i\to f\), alors \( f\) est continue.
\end{proposition}

%+++++++++++++++++++++++++++++++++++++++++++++++++++++++++++++++++++++++++++++++++++++++++++++++++++++++++++++++++++++++++++
\section{Ensembles nulle part denses}
%+++++++++++++++++++++++++++++++++++++++++++++++++++++++++++++++++++++++++++++++++++++++++++++++++++++++++++++++++++++++++++

Nous allons nous limiter au cas de \( \eR\), mais je crois que ça se généralise sans trop de peine aux espaces métriques, voire plus. Voir aussi la section~\ref{SecBDlaUrz} sur les espaces de Baire.

\begin{definition}
    Un ensemble est dit \defe{nulle part dense}{nulle part dense}\index{dense!nulle part} s'il n'est dense dans aucun intervalle.

    Un ensemble dans \( \eR\) est de \defe{première catégorie}{catégorie!ensemble de première} ou \defe{maigre}{maigre (ensemble)} s'il est une union dénombrable d'ensembles nulle part dense (c'est à dire d'ensembles denses sur aucun intervalle).
\end{definition}

\begin{theorem}[Baire\cite{BaireZied}]      \label{ThoQGalIO}
    Une réunion dénombrable d'ensembles nulle part denses est d'intérieur vide.
\end{theorem}
\index{Baire!théorème}
\index{théorème!Baire}

\begin{proof}
    Soient \( a\in S\) et \( \epsilon>0\). Nous allons trouver un élément dans \( B(a,\epsilon)\) qui n'est pas dans \( S\). Nous commençons par choisir \( x_1\in B(a,\epsilon)\) et \( r_1<\frac{ \epsilon }{2}\) tel que
    \begin{equation}
        B(x_1,r_1)\cap A_1=\emptyset.
    \end{equation}
    Ensuite nous choisissons \( x_2\in B(x_1,r_1)\) et \( r_2<\epsilon/4\) tel que \( B(x_2,r_2)\subset B(x_1,r_1)\) et \( B(x_2,r_2)\cap A_2=\emptyset\). Notons que \( B(x_2,r_2)\cap A_1=\emptyset\) aussi, par construction.

    Par récurrence nous construisons une suite d'éléments \( x_n\) et de rayons \( r_n<\epsilon/2^n\) tels que
    \begin{enumerate}
        \item
            \( B(x_n,r_n)\cap A_j=\emptyset\) pour tout \( j\leq n\),
        \item
            \( \overline{ B(x_n,r_n) }\subset B(x_{n-1},r_{r-1})\).
    \end{enumerate}
    Cette suite étant de Cauchy (parce que contenue dans des intervalles emboîtés de rayon décroissant vers zéro), elle converge\footnote{Par la proposition~\ref{PROPooTFVOooFoSHPg}} donc vers un point qui en particulier appartient à \( B(a,\epsilon)\). Mais la limite n'est dans aucun des \( A_n\) et donc pas dans \( S\).
\end{proof}

%+++++++++++++++++++++++++++++++++++++++++++++++++++++++++++++++++++++++++++++++++++++++++++++++++++++++++++++++++++++++++++
\section{Topologie des semi-normes}
%+++++++++++++++++++++++++++++++++++++++++++++++++++++++++++++++++++++++++++++++++++++++++++++++++++++++++++++++++++++++++++

Les principaux espaces topologiques construit avec des semi-normes seront les espaces de fonctions de la définition~\ref{DefFGGCooTYgmYf}. Nous verrons également la topologie \( *\)-faible sur \( \swD'(\Omega)\) en la définition~\ref{DefASmjVaT}.

\begin{definition}  \label{DefPNXlwmi}
    Si \( E\) est un espace vectoriel, une \defe{semi-norme}{semi-norme} sur \( E\) est une application \( p\colon E\to \eR\) telle que
    \begin{enumerate}
        \item
            \( p(x)\geq 0\),
        \item   \label{ItemSHnimhDii}
            \( p(\lambda x)=| \lambda |p(x)\)
        \item   \label{ItemSHnimhDiii}
            \( p(x+y)\leq p(x)+p(y)\).
    \end{enumerate}
\end{definition}

La seule différence avec une norme est simplement qu'une semi-norme peut s'annuler en des éléments non-nuls de l'espace.

\begin{lemma}[\cite{DRcmzcB}]
    Si \( p\) est une semi-norme nous avons
    \begin{equation}
        | p(x)-p(y) |\leq p(x-y).
    \end{equation}
\end{lemma}

\begin{proof}
    Nous avons d'une part \( p(x+h)\leq p(x)+p(h)\) et d'autre part \( p(x)\leq p(x+h)+p(-h)=p(x+h)+p(h)\). En isolant \( p(x+h)-p(x)\) dans chacune ce des deux inégalités,
    \begin{equation}
        -p(h)\leq p(x+h)-p(x)\leq p(h)
    \end{equation}
    ou encore
    \begin{equation}
        |p(x+h)-p(x)|\leq p(h)
    \end{equation}
    qui donne le résultat demandé en posant \( h=y-x\).
\end{proof}

Soit \( (p_i)_{i\in I}\) une famille de semi-normes sur \( E\). Nous construisons alors une topologie sur \( E\) de la façon suivante.

\begin{definition}[Topologie et semi-normes\cite{SOdaAdx,MUbDonp}]
    Pour tout \( J\) fini dans \( I\) nous définissons les \defe{boules ouvertes}{boule!avec semi-normes}
    \begin{equation}
        B_J(x,r)=\{ y\in E\tq p_j(y-x)<r\,\forall j\in J \}.
    \end{equation}
    La \defe{topologie}{topologie!et semi-normes} sur \( E\) donnée par la famille de semi-norme est définie en disant que \( \mO\subset E\) est ouvert si et seulement si chaque point de \( \mO\) est dans une boule contenue dans \( \mO\).
\end{definition}

\begin{proposition} \label{PropQPzGKVk}
    Une suite \( (x_n)\) dans \( E\) converge vers \( x\) au sens de la topologie des semi-normes si et seulement si pour tout \( i\in I\),
    \begin{equation}
        p_i(x-x_n)\to 0.
    \end{equation}
\end{proposition}

\begin{proof}
    Si la suite \( (x_n)\) converge\footnote{Définition~\ref{DefXSnbhZX}.} vers \( x\), alors pour tout ouvert \( \mO\) autour de \( x\), il existe un \( N\) tel que si \( n\geq N\), alors \( x_n\in\mO\). En particulier pour tout \( j\) et pour tout \( \epsilon>0\), il doit exister un \( n\geq N_j\) tel que \( x_n\in B_j(x,\epsilon)\).

    Voyons l'implication inverse. Soit \( \epsilon>0\). Pour tout \( i\in I\), il existe un \( N_i\) tel que \( n\geq N_i\) implique \( p_i(x-x_n)\leq \epsilon\). Si \( \mO\) est un ouvert, il doit contenir une boule du type \( B_J(x,r)\) pour un certain ensemble fini \( J\subset I\).

    En prenant \( N=\max\{ N_j\tq j\in J \}\), nous avons \( p_j(x-x_n)\leq \epsilon\) pour tout \( j\) et donc \( x_n\in B_J(x,r)\).
\end{proof}

La proposition suivante est un vulgaire plagiat de la proposition~\ref{PropQZRNpMn}.
\begin{proposition} \label{PropNGjQnqF}
    Soit \( f\colon \eR\to (E,p_i)_{i\in I}\) une application. Nous avons équivalence entre
    \begin{enumerate}
        \item   \label{ItemHNxGMpCi}
            la fonction \( f\) est continue en \( t_0\in \eR\),
        \item\label{ItemHNxGMpCii}
            si \( W\) est un voisinage ouvert de \( f(t_0)\) il existe un voisinage ouvert \( V\) de \( t_0\) (dans \( \eR\)) tel que \( f(V)\subset W\),
        \item\label{ItemHNxGMpCiii}
            pour tout \( i\in I\) et \( \epsilon>0\) il existe \( \delta>0\) tel que
            \begin{equation}
                f\big( B(t_0,\delta) \big)\subset B_i\big( f(t_0),\epsilon \big).
            \end{equation}
    \end{enumerate}
\end{proposition}

\begin{proof}
    L'équivalence~\ref{ItemHNxGMpCi} \( \Leftrightarrow\)~\ref{ItemHNxGMpCii} est la définition~\ref{DefOLNtrxB}.

    Prouvons~\ref{ItemHNxGMpCii} \( \Rightarrow\)~\ref{ItemHNxGMpCiii}. Soient \( i\in I\) et \( \epsilon>0\). Considérons la boule \( B_i\big( f(t_0),\epsilon \big)\), qui est un ouvert de \( E\) contenant \( f(t_0)\). Il existe donc un ouvert \( V\) autour de \( t_0\) tel que \( f(V)\subset B_i\big( f(t_0),\epsilon \big)\). En particulier \( V\) contient une boule \( B(t_0,\delta)\) et nous avons
    \begin{equation}
        f\big( B(t_0,\delta) \big)\subset f(V)\subset B_i\big( f(t_0),\epsilon \big).
    \end{equation}

    Prouvons~\ref{ItemHNxGMpCiii} \( \Rightarrow\)~\ref{ItemHNxGMpCii}. Soit \( W\) un ouvert autour de \( f(t_0)\). Il existe un \( i\in I\) et \( \epsilon>0\) tel que \( B_i\big( f(t_0),\epsilon \big)\subset W\). Nous avons alors un \( \delta>0\) tel que
    \begin{equation}
        f\big( B(t_0,\delta) \big)\subset B_i\big( f(t_0),\epsilon \big)\subset W.
    \end{equation}
\end{proof}

Lorsqu'on a un espace \( E\) muni d'une quantité dénombrable de semi-normes \( \{ p_k \}_{k\in I}\) nous définissons l'écart\footnote{Dans le cas de \( E=\swD(K)\), la première semi-norme est numérotée à zéro, donc il faudra poser \( d(\varphi_1,\varphi_2)\) avec \( p_{k-1}\) au lieu de \( p_k\).}
\begin{equation}        \label{EqAAghiUR}
    d(x,y)=\sup_{k\geq 1}\min\big\{  \frac{1}{ k },p_k(x-y) \big\}.
\end{equation}
Notons que cette écart est invariant par translation au sens où pour tout \( x,y,h\) dans \( E\) nous avons
\begin{equation}
    d(x+h,y+h)=\sup_{k\geq 1}\min\big\{ \frac{1}{ k },p_k(x-y) \big\}=d(x,y).
\end{equation}


\begin{proposition}     \label{PROPooMJEQooHtIyeX}
    Si \( X\) est un espace topologique dont la topologie est donnée par une famille dénombrable de semi-normes, alors il est métrisable.
\end{proposition}
%TODO : une preuve

\begin{proposition}[\cite{DRcmzcB}] \label{PropLOwUvCO}
    La topologie donnée par les boules
    \begin{equation}    \label{EqGHfYIlQ}
        B_k(a,r)=\{ x\in E\tq \forall\,k\leq \frac{1}{ r },p_k(x-a)<r\}
    \end{equation}
    est la même que celle «usuelle» donnée par les semi-normes. En disant «la même» nous entendons le fait que les ouverts sont les mêmes : \( A\) est ouvert pour une des deux topologies si et seulement s'il est ouvert pour l'autre.
\end{proposition}

\begin{proof}
    Pour cette démonstration nous allons préfixer par \( d\) les notions topologiques issues des boules \eqref{EqGHfYIlQ} et par \( P\) celle des semi-normes : \( P\)-continue, \( d\)-ouvert, etc.

    D'abord nous avons
    \begin{equation}    \label{EqRIURpQo}
        B(a,r)=\bigcap_{k\leq \frac{1}{ k }}B_k(a,r).
    \end{equation}
    Si \( \mO\) est un \( d\)-ouvert, il contient une \( d\)-boule autour de chacun de ses points. Or d'après la formule \eqref{EqRIURpQo}, une \( d\)-boule est une intersection \emph{finie} de \( P\)-ouverts et donc est un \( P\)-ouvert par définition. Donc \( \mO\) contient un \( P\)-ouvert autour de tous ses points et est donc \( P\)-ouvert.

    Inversement nous supposons que \( \mO\) est un \( P\)-ouvert. Commençons par prouver que les semi-normes \( p_k\) sont \( d\)-continues. En effet soient \( k\in \eN\), \( \epsilon\leq \frac{1}{ k }\) et \( x,y\in E\) tels que \( d(x,y)\leq \epsilon\); nous avons
    \begin{subequations}
        \begin{align}
            | p_k(y)-p_k(x) |&\leq p_k(x-y)\\
            &=\min\{ \frac{1}{ k },p_k(x-y) \}\\
            &\leq d(x,y)\\
            &\leq \epsilon.
        \end{align}
    \end{subequations}
    Montrons à présent que \( \mO\) est \( d\)-ouverte. Si \( a\in\mO\), il existe \( k\) et \( r\) tels que \( B_k(a,r)\subset\mO\). Soit \( x\in B_k(a,r)\). Montrons que si \( \epsilon\) est suffisamment petit, la \( d\)-boule \( B(x,\epsilon)\) est inclue à \( B_k(a,r)\). Pour cela prenons \( y\in B(x,\epsilon)\); nous avons
    \begin{equation}
        \big| p_k(a-x)-p_k(a-y) \big|\leq d(x,y)\leq \epsilon.
    \end{equation}
    Par conséquent le nombre \( p_k(a-y)\) est dans l'intervalle
    \begin{equation}
        p_k(a-x)\pm\epsilon
    \end{equation}
    et il suffit de prendre \( \epsilon<\frac{ r-p_k(a-x) }{2}\).
\end{proof}

%---------------------------------------------------------------------------------------------------------------------------
\subsection{Espace dual}
%---------------------------------------------------------------------------------------------------------------------------

Nous parlerons plus en détail d'espace dual d'un espace normé en la section~\ref{SECooKOJNooQVawFY}.

\begin{definition}  \label{DefHUelCDD}
    Soient \( F\) un espace métrique et \( E\) un espace topologique vectoriel. Une topologie possible\footnote{C'est, dans l'idée, celle qui sera choisie pour les espaces de distributions, voir la définition~\ref{DefASmjVaT}.} sur l'espace des applications linéaires \( \aL(E,F)\) est la \defe{topologie \( *\)-faible}{topologie!$*$-faible} qui est la topologie des semi-normes
    \begin{equation}
        p_v(T)=\| T(v) \|_F.
    \end{equation}
\end{definition}
C'est une famille de semi-normes indicées par les éléments de \( E\). Si \( E\) est un espace métrique, c'est cette topologie qui sera considérée sur son dual topologique\index{topologie!sur dual topologique} \( E'\) des applications continues \( E\to \eR\).

La proposition suivante indique qu'elle est un peu la topologie de la convergence ponctuelle.
\begin{proposition}
    Soient \( E\) un espace muni de la topologie des semi-normes \( \{ p_i \}_{i\in I}\) et \( F\) un espace métrique. Soient une suite \( (T_n)\) dans \( \aL(E,F)\) et \( T\in \aL(E,F)\). Nous avons \( T_n\stackrel{*}{\longrightarrow}T\) si et seulement si \( T_n(v)\stackrel{F}{\longrightarrow}T(v)\) pour tout \( v\in E\).
\end{proposition}

\begin{proof}
    Nous avons équivalence entre les lignes suivantes :
    \begin{subequations}
        \begin{align}
            T_n\stackrel{*}{\longrightarrow}T\\
            p_v(T_n-T)\to 0\,\forall v\in E &&\text{proposition~\ref{PropQPzGKVk}}\\
            \| T_n(v)-T(v) \|_E\to 0\,\forall v\in E\\
            T_n(v)\stackrel{E}{\longrightarrow}T(v).
        \end{align}
    \end{subequations}
\end{proof}

%---------------------------------------------------------------------------------------------------------------------------
\subsection{Espace \texorpdfstring{$ C^k(\eR,E')$}{C(R,E')}}
%---------------------------------------------------------------------------------------------------------------------------

Nous revenons à nos histoires de limites de la définition~\ref{DefXSnbhZX}.
\begin{proposition}[Unicité de la limite dans un dual topologique] \label{PropRBCiHbz}
    Soient \( E\) un espace métrique et \( E'\) son dual topologique muni de sa topologie de la définition~\ref{DefHUelCDD}. Il y a unicité de l'élément de \( E'\) vers lequel une fonction \( u\colon \eR\to E' \) peut converger.
\end{proposition}

\begin{proof}
    Soit \( T\) un élément vers lequel \( u_t\) converge lorsque \( t\to t_0\). Soient \( \epsilon>0\) et \( x\in E\). La boule \( B_x(T,\epsilon)\) de \( E'\) subordonnée à la norme \( p_x\) et centrée en \( T\) est un ouvert de \( E'\). Étant donné que \( u\) converge vers \( T\) il existe \( \delta>0\) tel que \( u_t\in B_x(T,\epsilon)\) dès que \( | t-t_0 |\leq \delta\). Nous avons donc, pour tout \( x\in E\), la limite (dans \( \eR\)) :
    \begin{equation}
        \lim_{t\to t_0} u_t(x)=T(x).
    \end{equation}
    Cela prouve que la convergence de \( u\) vers \( T\) implique l'existence pour tout \( x\) de la limite de \( u_t(x)\) dans \( \eR\). Si \( T'\) est un autre élément vers lequel \( u_t\) converge, nous avons par le même raisonnement que
    \begin{equation}
        \lim_{t\to t_0} u_t(x)=T'(x).
    \end{equation}
    Par unicité de la limite dans \( \eR\)
    %TODO : prouver ça et mettre une référence.
    nous devons alors avoir \( T(x)=T'(x)\) pour tout \( x\), c'est à dire \( T=T'\).
\end{proof}

\begin{proposition} \label{PropVKSNflB}
    Soit \( u\colon \eR\to E'\) une fonction continue. Alors
    \begin{enumerate}
        \item   \label{ItemLSJjfZdi}
            pour tout \( x\in E\) la fonction \( t\mapsto u_t(x)\) est continue,
        \item\label{ItemLSJjfZdii}
            pour tout \( x\in E\) nous avons la limite dans \( \eR\)
            \begin{equation}    \label{EqWKdFPVO}
                \lim_{t\to t_0} u_t(x)=u_{t_0}(x),
            \end{equation}
        \item\label{ItemLSJjfZdiii}
            nous avons la limite dans \( E'\)
            \begin{equation}
                \lim_{t\to t_0} u_t=u_{t_0}.
            \end{equation}
    \end{enumerate}
\end{proposition}

\begin{proof}
    Soient \( x\in E\) et \( \epsilon> 0\). Par la proposition~\ref{PropNGjQnqF} la continuité de \( u\) donne un \( \delta>0\) tel que
    \begin{equation}
        u_{B(t_0,\delta)}\subset B_x(u_{t_0},\epsilon).
    \end{equation}
    C'est à dire que si \( | t-t_0 |\leq \delta\) nous avons
    \begin{equation}
        \big| u_{t_0}(x)-u_t(x) \big|<\epsilon,
    \end{equation}
    ce qui signifie bien que la fonction \( t\mapsto u_t(x)\) est continue en tant que fonction \( \eR\to \eR\). Cela est le point~\ref{ItemLSJjfZdi}. Le théorème de limite et continuité dans \( \eR\) nous donne immédiatement la limite \eqref{EqWKdFPVO}.
    % TODO : aussi fou que cela puisse paraître, le théorème limite et continuité dans R n'est pas prouvé ni énoncé. Le faire et lié vers ici.

    Nous passons à la preuve du point~\ref{ItemLSJjfZdiii}. Soit \( \mO\) un ouvert de \( E'\) contenant \( u_{t_0}\). Il existe donc un \( i\in I\) et \( \epsilon>0\) tel que \( B_i(u_{t_0},\epsilon)\subset \mO\). Étant donné que \( u\) est continue, il existe \( \delta>0\) tel que
    \begin{equation}
        u_{B(t_0,\delta)}\subset B_i(u_{t_0},\epsilon)\subset \mO.
    \end{equation}
    Cela signifie bien que
    \begin{equation}
        | t-t_0 |\leq \delta\Rightarrow u_t\in\mO,
    \end{equation}
    c'est à dire que nous avons la limite \( \lim_{t\to t_0} u_t=u_{t_0}\) dans \( E'\). Pour dire cela nous avons utilisé la définition~\ref{DefYNVoWBx} de la limite et le résultat d'unicité~\ref{PropRBCiHbz}.
\end{proof}

\begin{definition}  \label{DefDZsypWu}
    Si nous avons une application \( u\colon \eR\to E'\) nous considérons sa \defe{dérivée}{dérivée!fonction à valeurs dans $E'$} donnée par la limite
    \begin{equation}
        u'_{t_0}=\lim_{t\to t_0} \frac{ u_t-u_{t_0} }{ t-t_0 }.
    \end{equation}
    Cela est un nouvel élément de \( E'\) (pour peu que la limite existe). La fonction \( u'\colon \eR\to E'\) ainsi définie peut être continue ou non. Cela nous permet de définir les espaces \( C^k(\eR,E')\) et \( C^{\infty}(\eR,E')\).
\end{definition}
Une des principales utilisations que nous ferons de ces espaces seront les espaces de fonctions à valeurs dans les distributions tempérées dont nous parlerons dans la section~\ref{SecTEgDVWO}.

%+++++++++++++++++++++++++++++++++++++++++++++++++++++++++++++++++++++++++++++++++++++++++++++++++++++++++++++++++++++++++++
\section{Espaces de Baire}
%+++++++++++++++++++++++++++++++++++++++++++++++++++++++++++++++++++++++++++++++++++++++++++++++++++++++++++++++++++++++++++
\label{SecBDlaUrz}

\begin{definition}
    Un \defe{espace de Baire}{espace!de Baire}\index{Baire!espace} est un espace topologique dans lequel toute intersection dénombrable d'ouverts denses est dense.
\end{definition}

\begin{theorem}[Théorème de Baire\cite{SIdTHwW}]    \label{ThoBBIljNM}
    Les espaces suivants sont de Baire :
    \begin{enumerate}
        \item
            les espaces topologiques localement compacts,
        \item
            les espaces métriques complets (donc ceux de Banach en particulier),
        \item
            tout ouvert d'un espace de Baire.
    \end{enumerate}
\end{theorem}
\index{théorème!de Baire}
\index{Baire!théorème}
%TODO : une preuve c'est sans doute bien, et ça a l'air d'être pas trop dur et donné sur Wikipédia.

\begin{proof}
    \begin{subproof}
    \item[Espaces topologiques localement compacts]
        \item[Espaces métriques complets]
            Soit \( (E,d)\) un espace métrique complet. Soient \( V\) un ouvert quelconque de \( E\) et \( U_n\) une suite d'ouverts denses. Le but est de prouver que l'ensemble \( \bigcap_{n\in \eN}U_n\) intersecte \( V\). Vu que \( V\) est ouvert dans un espace métrique, il contient une boule ouverte et donc une boule fermée \( B_0\) de rayon strictement positif. L'ensemble \( U_1\) est dense et intersecte donc un ouvert contenu dans \( B_0\). L'intersection est un ouvert qui contient alors une boule fermée \( B_1\) de rayon strictement positif. Continuant ainsi nous construisons une suite de fermés emboités \( B_n\) telle que
            \begin{equation}
                \bigcap_{n\in \eN}U_n\cap V
            \end{equation}
            contient l'intersection des \( B_n\). Par le théorème~\ref{ThoCQAcZxX} des fermés emboîtés (que nous utilisons parce que \( E\) est métrique et complet), cette intersection est non vide.
        \item[Ouvert d'un espace de Baire]
    \end{subproof}
\end{proof}

Une des applications du théorème de Baire est le théorème de Banach-Steinhaus~\ref{ThoPFBMHBN}.

%+++++++++++++++++++++++++++++++++++++++++++++++++++++++++++++++++++++++++++++++++++++++++++++++++++++++++++++++++++++++++++
\section{Sommes de familles infinies}
%+++++++++++++++++++++++++++++++++++++++++++++++++++++++++++++++++++++++++++++++++++++++++++++++++++++++++++++++++++++++++++
\label{SECooHHDXooUgLhHR}

%---------------------------------------------------------------------------------------------------------------------------
\subsection{Convergence commutative}
%---------------------------------------------------------------------------------------------------------------------------

\begin{definition}
    Soit \( x_k\) une suite dans un espace vectoriel normé \( E\). Nous disons que la suite \defe{converge commutativement}{convergence!commutative} vers \( x\in E\) si \( \lim_{n\to \infty}\| x_n-x \| =0\) et si pour toute bijection \( \tau\colon \eN\to \eN\) nous avons aussi
    \begin{equation}
        \lim_{n\to \infty} \| x_{\tau(k)}-x \|=0.
    \end{equation}
    La notion de convergence commutative est surtout intéressante pour les séries. La somme
    \begin{equation}
        \sum_{k=0}^{\infty}x_k
    \end{equation}
    converge commutativement vers \( x\) si \( \lim_{N\to \infty} \| x-\sum_{k=0}^Nx_k \|=0\) et si pour toute bijection \( \tau\colon \eN\to \eN\) nous avons
    \begin{equation}
        \lim_{N\to \infty} \| x-\sum_{k=0}^Nx_{\tau(k)} \|=0.
    \end{equation}
\end{definition}

Nous démontrons maintenant qu'une série converge commutativement si et seulement si elle converge absolument.

Pour le sens inverse, nous avons la proposition suivante.
\begin{proposition}
    Soit \( \sum_{k=0}^{\infty}a_k\) une série réelle qui converge mais qui ne converge pas absolument. Alors pour tout \( b\in \eR\), il existe une bijection \( \tau\colon \eN\to \eN\) telle que \( \sum_{i=0}^{\infty}a_{\tau(i)}=b\).
\end{proposition}
Pour une preuve, voir \href{http://gilles.dubois10.free.fr/analyse_reelle/seriescomconv.html}{chez Gilles Dubois}.

\begin{proposition} \label{PopriXWvIY}
    Soit \( (a_i)_{i\in \eN}\) une suite absolument convergente dans \( \eC\). Alors elle converge commutativement.
\end{proposition}

\begin{proof}
    Soit \( \epsilon>0\). Nous posons \( \sum_{i=0}^\infty a_i=a\) et nous considérons \( N\) tel que
    \begin{equation}
        | \sum_{i=0}^Na_i-a |<\epsilon.
    \end{equation}
    Étant donné que la série des \( | a_i |\) converge, il existe \( N_1\) tel que pour tout \( p,q>N_1\) nous ayons \( \sum_{i=p}^q| a_i |<\epsilon\). Nous considérons maintenant une bijection \( \tau\colon \eN\to \eN \). Prouvons que la série \( \sum_{i=0}^{\infty}| a_{\tau(i)} |\) converge. Nous choisissons \( M\) de telle sorte que pour tout \( n>M\), \( \tau(n)>N_1\); alors si \( p,q>M\) nous avons
    \begin{equation}
        \sum_{i=p}^q| a_{\tau(i)} |<\epsilon.
    \end{equation}
    Par conséquent la somme de la suite \( (a_{\tau(i)})\) converge. Nous devons montrer à présent qu'elle converge vers la même limite que la somme «usuelle» \( \lim_{N\to \infty} \sum_{i=0}^Na_i\).

    Soit \( n>\max\{ M,N \}\). Alors
    \begin{equation}
        \sum_{k=0}^na_{\tau(k)}-\sum_{k=0}^na_k=\sum_{k=0}^Ma_{\tau(k)}-\sum_{k=0}^Na_k+\underbrace{\sum_{M+1}^na_{\tau(k)}}_{<\epsilon}-\underbrace{\sum_{k=N+1}^na_k}_{<\epsilon}.
    \end{equation}
    Par construction les deux derniers termes sont plus petits que \( \epsilon\) parce que \( M\) et \( N\) sont les constantes de Cauchy pour les séries \( \sum a_{\tau(i)}\) et \( \sum a_i\). Afin de traiter les deux premiers termes, quitte à redéfinir \( M\), nous supposons que \( \{ 1,\ldots, N \}\subset \tau\{ 1,\ldots, M \}\); par conséquent tous les \( a_i\) avec \( i<N\) sont atteints par les \( a_{\tau(i)}\) avec \( i<M\). Dans ce cas, les termes qui restent dans la différence
    \begin{equation}
        \sum_{k=0}a_{\tau(k)}-\sum_{k=0}^Na_k
    \end{equation}
    sont des \( a_k\) avec \( k>N\). Cette différence est donc en valeur absolue plus petite que \( \epsilon\), et nous avons en fin de compte que
    \begin{equation}
        \left| \sum_{k=0}^na_{\tau(k)}-\sum_{k=0}^na_k \right| <\epsilon.
    \end{equation}
\end{proof}

\begin{proposition}     \label{PropyFJXpr}
    Soit \( \sum_{i=0}^{\infty}a_i\) une série qui converge mais qui ne converge pas absolument. Pour tout \( b\in \eR\), il existe une bijection \( \tau\colon \eN\to \eN\) telle que \( \sum_{i=}^{\infty}a_{\tau(i)}=b\).
\end{proposition}

Les propositions~\ref{PopriXWvIY} et~\ref{PropyFJXpr} disent entre autres qu'une série dans \( \eC\) est commutativement sommable si et seulement si elle est absolument sommable.

Soit \( (a_i)_{i\in I}\) une famille de nombres complexes indexée par un ensemble \( I\) quelconque. Nous allons nous intéresser à la somme \( \sum_{i\in I}a_i\).


Soit \( \{ a_i \}_{i\in I}\) des nombres positifs. Nous définissons la somme
\begin{equation}
    \sum_{i\in I}a_i=\sup_{ J\text{ fini}}\sum_{j\in J}a_j.
\end{equation}
Notons que cela est une définition qui ne fonctionne bien que pour les sommes de nombres positifs. Si \( a_i=(-1)^i\), alors selon la définition nous aurions \( \sum_i(-1)^i=\infty\). Nous ne voulons évidemment pas un tel résultat.

Dans le cas de familles de nombres réels positifs, nous avons une première définition de la somme.
\begin{definition}  \label{DefHYgkkA}
Soit \( (a_i)_{i\in I}\) une famille de nombres réels positifs indexés par un ensemble quelconque \( I\). Nous définissons
\begin{equation}
    \sum_{i\in I}a_i=\sup_{ J\text{ fini dans } I}\sum_{j\in J}a_j.
\end{equation}
\end{definition}

\begin{definition}  \label{DefIkoheE}
    Si \( \{ v_i \}_{i\in I}\) est une famille de vecteurs dans un espace vectoriel normé indexée par un ensemble quelconque \( I\). Nous disons que cette famille est \defe{sommable}{famille!sommable} de somme \( v\) si pour tout \( \epsilon>0\), il existe un \( J_0\) fini dans \( I\) tel que pour tout ensemble fini \( K\) tel que \( J_0\subset K\) nous avons
    \begin{equation}
        \| \sum_{j\in K}v_j-v \|<\epsilon.
    \end{equation}
\end{definition}
Notons que cette définition implique la convergence commutative.

\begin{example}
    La suite \( a_i=(-1)^i\) n'est pas sommable parce que quel que soit \( J_0\) fini dans \( \eN\), nous pouvons trouver \( J\) fini contenant \( J_0\) tel que \( \sum_{j\in J}(-1)^j>10\). Pour cela il suffit d'ajouter à \( J_0\) suffisamment de termes pairs. De la même façon en ajoutant des termes impairs, on peut obtenir \( \sum_{j\in J'}(-1)^i<-10\).
\end{example}

\begin{example}
    De temps en temps, la somme peut sortir d'un espace. Si nous considérons l'espace des polynômes \( \mathopen[ 0 , 1 \mathclose]\to \eR\) muni de la norme uniforme, la somme de l'ensemble
    \begin{equation}
        \{ 1,-1,\pm\frac{ x^n }{ n! } \}_{n\in \eN}
    \end{equation}
    est zéro.

    Par contre la somme de l'ensemble \( \{ 1,\frac{ x^n }{ n! } \}_{n\in \eN}\) est l'exponentielle qui n'est pas un polynôme.
\end{example}

\begin{example}
    Au sens de la définition~\ref{DefIkoheE} la famille
    \begin{equation}
        \frac{ (-1)^n }{ n }
    \end{equation}
    n'est pas sommable. En effet la somme des termes pairs est \( \infty\) alors que la somme des termes impairs est \( -\infty\). Quel que soit \( J_0\in \eN\), nous pouvons concocter, en ajoutant des termes pairs, un \( J\) avec \( J_0\subset J\) tel que \( \sum_{j\in J}(-1)^j/j\) soit arbitrairement grand. En ajoutant des termes négatifs, nous pouvons également rendre \( \sum_{j\in J}(-1)^j/j\) arbitrairement petit.
\end{example}

\begin{proposition} \label{PropVQCooYiWTs}
    Si \( (a_{ij})\) est une famille de nombres positifs indexés par \( \eN\times \eN\) alors
    \begin{equation}
        \sum_{(i,j)\in \eN^2}a_{ij}=\sum_{i=1}^{\infty}\Big( \sum_{j=1}^{\infty}a_{ij} \Big)
    \end{equation}
    où la somme de gauche est celle de la définition~\ref{DefHYgkkA}.
\end{proposition}
%TODO : cette proposition peut être vue comme une application de Fubini pour la mesure de comptage. Le faire et référentier ici.

\begin{proof}
    Nous considérons \( J_{m,n}=\{ 0,\ldots, m \}\times \{ 0,\ldots, n \}\) et nous avons pour tout \( m\) et \( n\) :
    \begin{equation}
        \sum_{(i,j)\in \eN^2}a_{ij}\geq \sum_{(i,j)\in J_{m,n}}a_{ij}=\sum_{i=1}^m\Big( \sum_{j=1}^na_{ij} \Big).
    \end{equation}
    Si nous fixons \( m\) et que nous prenons la limite \( n\to \infty\) (qui commute avec la somme finie sur \( i\)) nous trouvons
    \begin{equation}
        \sum_{(i,j)\in \eN^2}a_{ij}\geq =\sum_{i=1}^m\Big( \sum_{j=1}^{\infty}a_{ij} \Big).
    \end{equation}
    Cela étant valable pour tout \( m\), c'est encore valable à la limite \( m\to \infty\) et donc
    \begin{equation}
        \sum_{(i,j)\in \eN^2}a_{ij}\geq \sum_{i=1}^{\infty}\Big( \sum_{j=1}^{\infty}a_{ij} \Big).
    \end{equation}

    Pour l'inégalité inverse, il faut remarquer que si \( J\) est fini dans \( \eN^2\), il est forcément contenu dans \( J_{m,n}\) pour \( m\) et \( n\) assez grand. Alors
    \begin{equation}
        \sum_{(i,j)\in J}a_{ij}\leq \sum_{(i,j)\in J_{m,n}}a_{ij}=\sum_{i=1}^m\sum_{j=1}^na_{ij}\leq \sum_{i=1}^{\infty}\Big( \sum_{j=1}^{\infty}a_{ij} \Big).
    \end{equation}
    Cette inégalité étant valable pour tout ensemble fini \( J\subset \eN^2\), elle reste valable pour le supremum.
\end{proof}

La définition générale de la somme~\ref{DefIkoheE} est compatible avec la définition usuelle dans les cas où cette dernière s'applique.
\begin{proposition}[commutative sommabilité]\label{PropoWHdjw}
    Soit \( I\) un ensemble dénombrable et une bijection \( \tau\colon \eN\to I\). Soit \( (a_i)_{i\in I}\) une famille dans un espace vectoriel normé. Alors
    \begin{equation}
        \sum_{k=0}^{\infty}a_{\tau(k)}=\sum_{i\in I}a_i
    \end{equation}
    dès que le membre de droite existe. Le membre de gauche est définit par la limite usuelle.
\end{proposition}

\begin{proof}
    Nous posons \( a=\sum_{i\in I}a_i\). Soit \( \epsilon>0\) et \( J_0\) comme dans la définition. Nous choisissons
    \begin{equation}
        N>\max_{j\in J_0}\{ \tau^{-1}(j) \}.
    \end{equation}
    En tant que sommes sur des ensembles finis, nous avons l'égalité
    \begin{equation}
        \sum_{k=0}^Na_{\tau(k)}=\sum_{j\in J_0}a_j
    \end{equation}
    où \( J\) est un sous-ensemble de \( I\) contenant \( J_0\). Soit \( J\) fini dans \( I\) tel que \( J_0\subset J\). Nous avons alors
    \begin{equation}
        \| \sum_{k=0}^Na_{\tau(k)}-a \|=\| \sum_{j\in J}a_j-a \|<\epsilon.
    \end{equation}
    Nous avons prouvé que pour tout \( \epsilon\), il existe \( N\) tel que \( n>N\) implique \( \| \sum_{k=0}^na_{\tau(k)}-a\| <\epsilon\).
\end{proof}

\begin{corollary}
    Nous pouvons permuter une somme dénombrable et une fonction linéaire continue. C'est à dire que si \( f\) est une fonction linéaire continue sur l'espace vectoriel normé \( E\) et \( (a_i)_{i\in I}\) une famille sommable dans \( E\) alors
    \begin{equation}
        f\left( \sum_{i\in I}a_i \right)=\sum_{i\in I}f(a_i).
    \end{equation}
\end{corollary}

\begin{proof}
    En utilisant une bijection \( \tau\) entre \( I\) et \( \eN\) avec la proposition~\ref{PropoWHdjw} ainsi que le résultat connu à propos des sommes sur \( \eN\), nous avons
    \begin{subequations}
        \begin{align}
            f\left( \sum_{i\in I}a_i \right)&=f\left( \sum_{k=0}^{\infty}a_{\tau(k)} \right)\\
            &=\sum_{k=0}^{\infty}f(a_{\tau(k)}) \label{SUBEQooCVUTooPmnHER}\\
            &=\sum_{i\in I}f(a_i).
        \end{align}
    \end{subequations}
    Notons que le passage à \eqref{SUBEQooCVUTooPmnHER} n'est pas du tout une trivialité à deux francs cinquante. Il s'agit d'écrire la somme comme la limite des sommes partielles, et de permuter \( f\) avec la limite en invoquant la continuité, puis de permuter \( f\) avec la somme partielle en invoquant sa linéarité.

    Ah, tiens et tant qu'on y est à dire qu'il y a des choses évidentes qui ne le sont pas, oui, il existe des applications linéaires non continues, voir le thème~\ref{THEMEooYCBUooEnFdUg}.
\end{proof}

La proposition suivante nous enseigne que les sommes infinies peuvent être manipulée de façon usuelle.
\begin{proposition} \label{PropMpBStL}
    Soit \( I\) un ensemble dénombrable. Soient \( (a_i)_{i\in I}\) et \( (b_i)_{i\in I}\), deux familles de réels positifs telles que \( a_i<b_i\) et telles que \( (b_i)\) est sommable. Alors \( (a_i)\) est sommable.

    Si \( (a_i)_{i\in I}\) est une famille de complexes telle que \( (| a_i |)\) est sommable, alors \( (a_i)\) est sommable.
\end{proposition}

\begin{proposition}[\cite{MonCerveau}]     \label{PROPooWLEDooJogXpQ}
    Soit un espace vectoriel normé \( E\) et une famille sommable\footnote{Définition~\ref{DefIkoheE}.} \( \{ v_i \}_{i\in I}\) d'éléments de \( E\). Soit \( f\colon E\to \eC\) une application sur laquelle nous supposons
    \begin{enumerate}
        \item
            \( f\) est linéaire et continue;
        \item
            la partie \( \{ f(v_i)_{i\in I} \} \) est sommable.
    \end{enumerate}
    Alors nous pouvons permuter la somme et \( f\) :
    \begin{equation}        \label{EQooONHXooKqIEbY}
        f\big( \sum_{i\in I}v_i \big)=\sum_{i\in I}f(v_i).
    \end{equation}
\end{proposition}

\begin{proof}
    Soit \( \epsilon>0\); vu que les familles \( \{ v_i \}_{i\in I}\) et \( \{ f(v_i) \}_{i\in I}\) sont sommables, nous pouvons considérer les parties finies \( J_1\) et \( J_2\) de \( I\) telles que
    \begin{equation}
        \big\| \sum_{j\in J_1}v_j-\sum_{i\in I}v_i \big\|\leq \epsilon
    \end{equation}
    et
    \begin{equation}
        \big\| \sum_{j\in J_2}f(v_j)-\sum_{i\in I}f(v_i) \big\|\leq \epsilon
    \end{equation}
    Ensuite nous posons \( J=J_1\cup J_2\). Avec cela nous calculons un peu avec les majorations usuelles :
    \begin{equation}
        \| f(\sum_{i\in I}v_i) -\sum_{i\in I}f(v_i) \|\leq \| f(\sum_{i\in I}v_i)- f(\sum_{j\in J}v_j) \|+  \| f(\sum_{j\in J}v_j)-\sum_i\in If(v_i) \|.
    \end{equation}
    Le second terme est majoré par \( \epsilon\), tandis que le premier, en utilisant la linéarité de \( f\) possède la majoration
    \begin{equation}
        \| f(\sum_{i\in I}v_i)- f(\sum_{j\in J}v_j) \|=\| f(\sum_{i\in I}v_i-\sum_{j\in J}v_j) \|\leq \| f \| \| \sum_{i\in I}v_i- \sum_{j\in J}v_j\|\leq \epsilon\| f \|.
    \end{equation}
    Donc pour tout \( \epsilon>0\) nous avons
    \begin{equation}
        \| f(\sum_{i\in I}v_i) -\sum_{i\in I}f(v_i) \|\leq \epsilon(1+\| f \|).
    \end{equation}
    D'où l'égalité \eqref{EQooONHXooKqIEbY}.
\end{proof}


%+++++++++++++++++++++++++++++++++++++++++++++++++++++++++++++++++++++++++++++++++++++++++++++++++++++++++++++++++++++++++++
\section{Fonctions}		\label{Sect_fonctions}
%+++++++++++++++++++++++++++++++++++++++++++++++++++++++++++++++++++++++++++++++++++++++++++++++++++++++++++++++++++++++++++

Soient $(V,\| . \|_V)$ et $(W,\| . \|_W)$ deux espaces vectoriels normés, et une fonction $f$ de $V$ dans $W$. Il est maintenant facile de définir les notions de limites et de continuité pour de telles fonctions en copiant les définitions données pour les fonctions de $\eR$ dans $\eR$ en changeant simplement les valeurs absolues par les normes sur $V$ et $W$.

La caractérisation suivante est un recopiage de la définition~\ref{DefOLNtrxB} lorsque la topologie est donnée par des boules.
\begin{proposition}\label{PropHOCWooSzrMjl}
	Soit $f\colon V\to W$ une fonction de domaine \( \Domaine(f)\subset V\) et soit $a$ un point d'accumulation de $\Domaine(f)$.
    La fonction \( f\) admet une limite en $a\in V$ si et seulement s'il existe un élément $\ell\in W$ tel que pour tout $\varepsilon>0$, il existe un $\delta>0$ tel que pour tout $x\in \Domaine(f)$,
    \begin{equation}        \label{EqDefLimzxmasubV}
		0<\| x-a \|_V<\delta\,\Rightarrow\,\| f(x)-\ell \|_W<\varepsilon.
	\end{equation}
	Dans ce cas, nous écrivons $\lim_{x\to a} f(x)=\ell$ et nous disons que $\ell$ est la \defe{limite}{limite} de $f$ lorsque $x$ tend vers $a$.
\end{proposition}

\begin{remark}
    Le fait que nous limitions la formule \eqref{EqDefLimzxmasubV} aux \( x\) dans le domaine de \( f\) n'est pas anodin. Considérons la fonction \( f(x)=\sqrt{x^2-4}\), de domaine \( | x |\geq 2\). Nous avons
    \begin{equation}
        \lim_{x\to 2} \sqrt{x^2-4}=0.
    \end{equation}
    Nous ne pouvons pas dire que cette limite n'existe pas en justifiant que la limite à gauche n'existe pas. Les points \( x<2\) sont hors du domaine de \( f\) et ne comptent dons pas dans l'appréciation de l'existence de la limite.

    Vous verrez plus tard que ceci provient de la \wikipedia{fr}{Topologie_induite}{topologie induite} de \( \eR\) sur l'ensemble \( \mathopen[ 2 , \infty [\).
\end{remark}